Mori–Zwanzig projection formalism
How irreversible, dissipative dynamics emerges from reversible mechanics
The Mori–Zwanzig (MZ) projection formalism is a systematic method for deriving effective equations of motion for a small number of physical variables from the dynamics of a much larger system.

Its fundamental result is that when some degrees of freedom are eliminated from a Hamiltonian system, the remaining variables generally obey equations containing three contributions:
\[
\boxed{
\underbrace{\dot A(t)}_{\text{reduced evolution}}
=
\underbrace{\text{instantaneous dynamics}}_{\text{reversible or mean force}}
+
\underbrace{\text{memory}}_{\text{history dependence}}
+
\underbrace{\text{fluctuating force}}_{\text{unresolved degrees of freedom}}.
}
\]

The important point is that this decomposition can be derived exactly, without initially assuming friction, randomness, or irreversible dynamics.

The formalism was developed by Robert Zwanzig and Hazime Mori, particularly during the 1960s, and remains central to nonequilibrium statistical mechanics, generalized Langevin equations, transport theory, and modern coarse-graining methods. 
Oxford Academic+1


It is especially relevant to our preceding discussion of metriplectic mechanics.

Metriplectic dynamics begins with an effective description combining reversible Hamiltonian evolution and entropy-producing dissipation.

Mori–Zwanzig addresses a more microscopic question:

Starting with an ordinary conservative Hamiltonian system, what happens to its equations of motion when we retain only some degrees of freedom?

We shall see that the immediate answer is not simply a dissipative force. The more fundamental result is nonlocality in time, or memory.
1. The physical problem: eliminating degrees of freedom
Consider an isolated Hamiltonian system with coordinates and momenta
\[
z=(q_1,\ldots,q_N,p_1,\ldots,p_N),
\]

and Hamiltonian
\[
H(z).
\]

Its microscopic evolution is
\[
\dot z^a=J^{ab}\partial_bH.
\]

Suppose that we are interested only in a small set of observables,
\[
A(z)=\big(A_1(z),\ldots,A_m(z)\big),
\qquad m\ll 2N.
\]

For example, these may be the position and momentum of a massive particle immersed in a fluid. The remaining variables describe the molecules of the fluid.

The exact equations are
\[
\dot A_i=\{A_i,H\}.
\]

But the right-hand side generally depends on microscopic variables not contained in \(A\). Therefore, the equations are not closed.

One might try to replace the unknown microscopic forces by averages. Unfortunately, this loses something essential.

The eliminated degrees of freedom retain information about their previous interactions with the observed subsystem. Later, that information can influence the subsystem again.

Therefore, the effective force at time \(t\) generally depends not just on \(A(t)\), but on its past history,
\[
A(s),\qquad 0\leq s\leq t.
\]

This is the origin of memory.

Furthermore, different microscopic initial conditions may have identical values of \(A(0)\), but subsequently produce different trajectories \(A(t)\).

This is the origin of the fluctuating force.

Neither effect requires fundamental stochasticity.
2. From Hamiltonian mechanics to the Liouville operator
Let \(\Phi_t\) denote the Hamiltonian flow in phase space:
\[
z(t)=\Phi_t(z_0).
\]

Instead of evolving phase-space points, we can evolve observables. For any function \(F(z)\),
\[
F(t,z_0)=F(\Phi_t(z_0)).
\]

Introduce the Liouville operator
\[
\boxed{
LF=\{F,H\}.
}
\]

For canonical coordinates,
\[
L
=
\sum_{i=1}^N
\left(
\frac{\partial H}{\partial p_i}
\frac{\partial}{\partial q_i}
-
\frac{\partial H}{\partial q_i}
\frac{\partial}{\partial p_i}
\right).
\]

Then
\[
\frac{dF}{dt}=LF
\]

along a trajectory, and the formal solution is
\[
\boxed{F(t)=e^{tL}F(0).}
\]

The operator \(e^{tL}\) is the Koopman evolution operator acting on observables.

Notice the distinction: \(L\) is a linear operator on functions even when the underlying Hamiltonian equations are nonlinear.

This linearity is essential to the Mori–Zwanzig construction.
3. Introducing the projection operator
We now divide observables into two classes:
- Relevant observables, expressible in terms of the retained variables \(A\).
- Irrelevant observables, containing additional information about the eliminated degrees of freedom.
Introduce a projection operator \(P\) and its complement \(Q\):
\[
\boxed{
P^2=P,\qquad Q=I-P,\qquad PQ=QP=0.
}
\]

The terminology projection deserves care.

This is not, in general, a projection of phase-space points onto a smaller phase space. It is a projection in a function space of observables.

There are two important versions.
A. Mori's linear projection
Suppose the equilibrium distribution is \(\rho_{\rm eq}(z)\), invariant under Hamiltonian evolution.

Define an inner product on observables:
\[
(F,G)=\int F^*(z)G(z)\rho_{\rm eq}(z)\,dz.
\]

For one centered observable \(A\), Mori's projector can be written
\[
\boxed{
PF=\frac{(A,F)}{(A,A)}A.
}
\]

This selects the component of \(F\) parallel to \(A\).

For several relevant observables, the construction uses their covariance matrix.

Mori's projector is particularly useful for deriving linear generalized Langevin equations and equilibrium fluctuation–dissipation relations.
B. Zwanzig's conditional expectation
Zwanzig's projection is more general:
\[
\boxed{
(PF)(z)=\mathbb E[F\mid A=A(z)].
}
\]

It averages over microscopic states compatible with a given value of the relevant variables.

For example, if \(A=(q,p)\) denotes the position and momentum of one particle, then \(PF\) averages over the remaining microscopic coordinates and momenta while holding \(q,p\) fixed.

This projection is linear as an operator on functions but generally produces nonlinear functions of \(A\).

It therefore permits nonlinear reduced evolution equations.

Neither construction requires a Riemannian metric on the physical configuration space. In Mori's case, the inner product is a statistical structure on observables.
4. Derivation of the exact Mori–Zwanzig equation
We begin with
\[
\frac{d}{dt}e^{tL}A=e^{tL}LA.
\]

Insert \(I=P+Q\):
\[
\frac{dA(t)}{dt}
=
e^{tL}PLA+e^{tL}QLA.
\]

The first term represents the projected instantaneous dynamics.

The second still contains unresolved information.

The crucial step is the Dyson–Duhamel identity:
\[
\boxed{
e^{tL}
=
e^{tQL}
+
\int_0^t
e^{(t-s)L}PL\,e^{sQL}\,ds.
}
\]

Applying this identity to \(QLA\) gives
\[
\begin{aligned}
\frac{dA(t)}{dt}
={}&e^{tL}PLA\\
&+\int_0^t e^{(t-s)L}
PL\,e^{sQL}QLA\,ds\\
&+e^{tQL}QLA.
\end{aligned}
\]

Thus we obtain the exact Mori–Zwanzig identity:
\[
\boxed{
\begin{aligned}
\dot A(t)={}&
\underbrace{e^{tL}PLA}_{\text{instantaneous term}}\\
&+
\underbrace{\int_0^t
e^{(t-s)L}PL\,e^{sQL}QLA\,ds}_{\text{memory term}}\\
&+
\underbrace{e^{tQL}QLA}_{\text{fluctuating term}}.
\end{aligned}
}
\]

No approximation has been made in this algebraic derivation, assuming the operators and their evolutions are sufficiently well-defined. The same identity applies formally to nonlinear dynamical systems. 
PMC+1


Define the fluctuating force
\[
R(t)=e^{tQL}QLA.
\]

It satisfies the auxiliary equation
\[
\dot R(t)=QLR(t),
\qquad R(0)=QLA.
\]

This is called the orthogonal dynamics equation.

Its solution is generally not ordinary Hamiltonian motion restricted to the eliminated variables. Rather, it is evolution under the operator \(QL\) in the space of observables.

For Zwanzig's conditional-expectation projector, the exact equation can be expressed as
\[
\boxed{
\dot A(t)=
b(A(t))
+\int_0^t K(A(t-s),s)\,ds
+R(t,z_0).
}
\]

Here \(b\) is the projected instantaneous drift, \(K\) is a generally nonlinear memory kernel, and \(R\) depends on the unresolved microscopic initial conditions.

This is a generalized Langevin equation, although \(R\) is not necessarily a stochastic process.
5. Mori's memory kernel and fluctuation–dissipation relation
Mori's linear projection gives a particularly instructive result.

For a single centered, real observable \(A\), consider
\[
PF=\frac{\langle AF\rangle}{\langle A^2\rangle}A.
\]

For an equilibrium Hamiltonian system, the Liouville operator is antisymmetric with respect to the equilibrium inner product:
\[
\langle F\,LG\rangle=-\langle (LF)\,G\rangle.
\]

Using this property, the exact projected equation becomes
\[
\boxed{
\dot A(t)
=
\Omega A(t)
-\int_0^t K(s)A(t-s)\,ds
+R(t),
}
\]

with
\[
\Omega=\frac{\langle ALA\rangle}{\langle A^2\rangle},
\]

and
\[
\boxed{
K(t)=\frac{\langle R(t)R(0)\rangle}
{\langle A^2\rangle}.
}
\]

For one real observable in an invariant equilibrium ensemble, \(\Omega=0\); for a collection of variables, it becomes a matrix describing reversible couplings.

The last equation is especially significant.

The memory kernel is determined by the autocorrelation of the fluctuating force.

The fluctuating and dissipative contributions therefore are not independent phenomenological additions.

They arise together from the same eliminated degrees of freedom.

This is one form of the second fluctuation–dissipation relation. 
Oxford Academic+1

6. Explicit Hamiltonian example: a particle coupled to an oscillator bath
Let us derive a generalized Langevin equation by explicitly eliminating microscopic degrees of freedom.

Consider a particle of mass \(M\), coordinate \(q\), and momentum \(p\), coupled to \(N\) harmonic oscillators.

Choose the Hamiltonian
\[
\boxed{
H=
\frac{p^2}{2M}+V(q)
+
\sum_{j=1}^N
\left[
\frac{p_j^2}{2m_j}
+
\frac{m_j\omega_j^2}{2}
\left(
x_j-\frac{c_jq}{m_j\omega_j^2}
\right)^2
\right].
}
\]

The shifted oscillator potentials include a counterterm that prevents an unwanted static renormalization of \(V(q)\).

The full system is Hamiltonian, reversible, and energy conserving.

Its equations are
\[
M\ddot q
=
-V'(q)
+\sum_j c_j
\left(x_j-\frac{c_jq}{m_j\omega_j^2}\right),
\]

and
\[
m_j\ddot x_j+m_j\omega_j^2x_j=c_jq.
\]

Eliminating the bath variables
The oscillator equations can be solved explicitly:
\[
\begin{aligned}
x_j(t)={}&x_j(0)\cos\omega_jt\\
&+\frac{p_j(0)}{m_j\omega_j}\sin\omega_jt\\
&+\frac{c_j}{m_j\omega_j}
\int_0^t\sin[\omega_j(t-s)]q(s)\,ds.
\end{aligned}
\]

Substituting these solutions into the equation for \(q\), and integrating by parts, produces
\[
\boxed{
M\ddot q(t)+V'(q(t))
+\int_0^t\Gamma(t-s)\dot q(s)\,ds
=\xi(t).
}
\]

The memory kernel is
\[
\boxed{
\Gamma(t)=
\sum_{j=1}^N
\frac{c_j^2}{m_j\omega_j^2}
\cos(\omega_jt).
}
\]

The fluctuating force is
\[
\boxed{
\begin{aligned}
\xi(t)=\sum_jc_j\bigg[
&\left(x_j(0)-\frac{c_jq(0)}{m_j\omega_j^2}\right)
\cos\omega_jt\\
&+\frac{p_j(0)}{m_j\omega_j}\sin\omega_jt
\bigg].
\end{aligned}
}
\]

This is an exact result.

It is closely related to Zwanzig's treatment of nonlinear generalized Langevin equations. 
courses.physics.ucsd.edu

What have we obtained?
We started from completely ordinary Hamiltonian mechanics.

After eliminating the oscillator variables, we obtained an equation involving:
1. The instantaneous conservative force \(-V'(q)\).
2. A history-dependent friction-like force.
3. A fluctuating force determined by microscopic initial conditions.
There is no stochastic assumption in this elimination. For specified initial conditions of every oscillator, \(\xi(t)\) is a definite function of time.

It becomes random only when we replace the unknown initial bath coordinates and momenta by a probability distribution.
Fluctuation–dissipation theorem
Suppose the bath is initially in conditional canonical equilibrium at temperature \(T\), with the particle's initial \(q,p\) held fixed.

Then
\[
\langle\xi(t)\rangle=0,
\]

and the force covariance is
\[
\boxed{
\langle\xi(t)\xi(s)\rangle
=
k_{\mathrm B}T\,\Gamma(t-s).
}
\]

This explicitly connects memory and fluctuations.

The same microscopic couplings \(c_j\) appear in both quantities. Neither can be chosen independently if the equilibrium microscopic model is to be respected.
7. How ordinary friction emerges
Suppose that the bath contains a very large number of modes whose correlations decay rapidly.

If the particle evolves much more slowly than the bath correlation time, we may approximate
\[
\int_0^t\Gamma(s)\dot q(t-s)\,ds
\approx
\dot q(t)\int_0^\infty\Gamma(s)\,ds.
\]

Define
\[
\gamma=\int_0^\infty\Gamma(s)\,ds,
\]

assuming the integral is well-defined in the relevant continuum-bath limit.

Then
\[
M\ddot q+\gamma\dot q+V'(q)=\xi(t).
\]

If the force correlations also become effectively delta-correlated, the result is the ordinary Langevin equation,
\[
\boxed{
M\ddot q+\gamma\dot q+V'(q)
=
\sqrt{2\gamma k_{\mathrm B}T}\,\eta(t),
}
\]

where \(\eta\) is normalized Gaussian white noise in the idealized limit.

Thus
\[
\text{Hamiltonian dynamics}
\longrightarrow
\text{generalized Langevin dynamics}
\longrightarrow
\text{Markovian Langevin dynamics}.
\]

The first step is an exact elimination of variables for the oscillator-bath model.

The second step requires physical approximations or suitable limiting procedures: decay of correlations, separation of time scales, and statistical assumptions about the bath.

There is an important qualification. For a finite oscillator bath, \(\Gamma(t)\) is a sum of cosines. The bath can return energy to the particle, and recurrences may occur.

Consequently, eliminating variables does not itself establish strict thermodynamic irreversibility.
8. Connection to metriplectic dynamics
We can now return to your original question.

Recall the metriplectic equation
\[
\dot z
=
J\nabla H+M\nabla S,
\]

with
\[
J\nabla S=0,\qquad M\nabla H=0.
\]

The symmetric positive semidefinite tensor \(M\) accounts for irreversible entropy-producing dynamics while conserving total energy.

Mori–Zwanzig, by contrast, starts from a microscopic Hamiltonian equation,
\[
\dot z=J\nabla H,
\]

and derives memory and fluctuating forces after projecting onto selected variables.

The conceptual relationship is:
Microscopic Hamiltonian system

Reversible evolution; all degrees of freedom retained

Mori–Zwanzig projection

Exact reduced equation

Instantaneous dynamics + memory + fluctuating force

Short-memory and statistical approximations

Effective dissipative dynamics

Friction, diffusion and transport equations

Thermodynamically consistent closure, when available

Metriplectic / GENERIC description

Energy-conserving reversible–irreversible dynamics with entropy production


The final step is not automatic.

An arbitrary Mori–Zwanzig memory kernel does not necessarily admit representation by a local, symmetric, positive semidefinite metric bracket.

Additional thermodynamic assumptions and suitable state variables are needed. In particular, for the isolated metriplectic formulation one must retain enough variables to represent both the total energy and entropy.

One advantage of Mori–Zwanzig is that it can, in principle, provide microscopic expressions for transport coefficients that a metriplectic formulation otherwise leaves phenomenological.

For example, friction and viscosity can often be related to time integrals of equilibrium correlation functions, leading to Green–Kubo formulas.
9. What Mori–Zwanzig does and does not explain
There are several points worth emphasizing.

First, the exact formalism is not intrinsically irreversible. The full microscopic dynamics remains Hamiltonian. Apparent irreversibility emerges in a reduced statistical description, typically after further approximations or appropriate limits.

Second, projection does not necessarily mean averaging over fast variables. It means selecting relevant observables. A useful exact projection can be constructed without separation of time scales, although short-memory approximations then become much less reliable.

Third, the fluctuating force is not fundamental noise. It is the effect of unresolved microscopic initial conditions. Its identification with a Gaussian or white-noise process requires further assumptions.

Fourth, exactness does not mean easy computability. Computing \(e^{tQL}\), the orthogonal dynamics, may be almost as difficult as solving the original microscopic system. The main practical challenge is finding controlled approximations for this term.

Recent mathematical work has also clarified that rigorous justification of the projection identities can depend on the operator domains and the choice of projector, particularly for nonlinear conditional-expectation projections. 
arXiv

10. References
These are particularly useful for a technical study.

1. R. Zwanzig, “Memory Effects in Irreversible Thermodynamics,” Physical Review 124, 983–992 (1961). Original paper. Note that a correction to part of its nonlinear transport derivation appeared in Physical Review A 5, 2680 (1972). 
Phys. Rev. A


2. H. Mori, “Transport, Collective Motion, and Brownian Motion,” Progress of Theoretical Physics 33, 423–455 (1965). Original paper. Fundamental for the linear projection and memory-kernel formulation.

3. R. Zwanzig, “Nonlinear Generalized Langevin Equations,” Journal of Statistical Physics 9, 215–220 (1973). Paper. Especially valuable for the oscillator-bath example.

4. R. Zwanzig, Nonequilibrium Statistical Mechanics, Oxford University Press (2001). Chapter 8, “Projection Operators,” pp. 143–168, gives a systematic presentation. Chapter

5. A. J. Chorin, O. H. Hald and R. Kupferman, “Optimal prediction and the Mori–Zwanzig representation of irreversible processes,” PNAS 97, 2968–2973 (2000). Open-access paper. Particularly interesting for its connection to mathematical model reduction.

6. E. Darve, J. Solomon and A. Kia, “Computing generalized Langevin equations and generalized Fokker–Planck equations,” PNAS 106, 10884–10889 (2009). Open-access paper. Discusses numerical treatment of projection and memory.

The deeper mechanical question
There is an important distinction between two statements:
\[
\text{Eliminating variables produces memory and fluctuations}
\]

and
\[
\text{Eliminating variables produces irreversible dissipation}.
\]
The first is an exact structural result of the Mori–Zwanzig formalism.

The second requires more: statistical assumptions, loss of accessible information, decay of correlations, or suitable limiting procedures.

This distinction is closely related to the longstanding problem of deriving irreversible thermodynamics from reversible classical mechanics.

For our discussion of metriplectic dynamics, a natural next technical problem would be to start with the Hamiltonian oscillator-bath model, derive its reduced Fokker–Planck equation in the Markovian limit, and then construct an explicit metriplectic formulation for the particle plus its thermal reservoir.

That would connect Hamiltonian mechanics, projection theory, statistical mechanics, and metriplectic geometry within one concrete calculation.
